\clearpage
\fancyhead[L]{\small IRIS.}
\SourceLeaf{479}
\begin{center}
\rule{0.55\linewidth}{0.4pt}\\[2pt]
Astronomische Nachrichten. Vol. XXVI. No. 611. pp. 173. 174. 16 Sept. 1847.\\[2pt]
\rule{0.55\linewidth}{0.4pt}
\end{center}
\begin{center}[Meridian-circle observations by \textsc{Gauss}.]\end{center}

\begin{center}
\small
\setlength{\tabcolsep}{2pt}
\begin{tabular}{@{}lrll@{}}
Iris: 1847 Aug. 21. & $9^{\rm u}\ 52^{\rm m}\ 54^{\rm s}1$ Göttingen mean time & $297^\circ\ 47'\ 20''8$ & $-13^\circ\ 42'\ 20''4$ \\
\end{tabular}
\end{center}

The timings at the individual wires did not agree particularly well because the object was hard to see through passing clouds. Immediately afterward the clouds spread still farther, so that the culmination of $\alpha$ Capricorni could not be observed.

\begin{center}
\small
\setlength{\tabcolsep}{4pt}
\begin{tabular}{lrrr}
Neptune 1847 Aug. 17. & $12^{\rm u}\ 22^{\rm m}\ 55^{\rm s}2$ & $331^\circ\ 27'\ 14''1$ & $-12^\circ\ 21'\ 5''0$ \\
\end{tabular}
\end{center}

